\documentclass[11pt]{article}

\title{Projective Incidence Geometry}
\author{UlamLib seed note}
\date{}

\newtheorem{definition}{Definition}[section]
\newtheorem{lemma}[definition]{Lemma}
\newtheorem{theorem}[definition]{Theorem}
\newtheorem{corollary}[definition]{Corollary}

\begin{document}
\maketitle

\section{Scope}

This note targets an incidence layer over existing projectivization objects.
The intended Lean destination is
\texttt{ULAM/Geometry/ProjectiveIncidence.lean}.

\section{Core definitions}

\begin{definition}
Three points in projective space are collinear if they lie on a common
projective line.
\end{definition}

\begin{definition}
Given two distinct points in projective space, the line through them is the
projectivization of the two-dimensional linear span of representatives of those
points.
\end{definition}

\section{Seed theorem cluster}

\begin{lemma}
Collinearity is symmetric in its arguments and every triple containing repeated
points is collinear.
\end{lemma}

\begin{theorem}[Two points determine a line]
Two distinct points of projective space lie on a unique projective line.
\end{theorem}

\begin{theorem}
Projectivizations of injective linear maps preserve collinearity and incidence.
\end{theorem}

\begin{corollary}
The standard projective space attached to a vector space satisfies the basic
incidence axioms of projective geometry.
\end{corollary}

\section{Formalization backlog}

\begin{enumerate}
\item Define collinearity.
\item Define projective lines using existing projective subspaces.
\item Prove uniqueness of the line through two distinct points.
\item Derive a projective-geometry axiom package.
\item Package incidence-preserving maps.
\end{enumerate}

\end{document}
